% preamble.tex
\documentclass[11pt, twoside]{book}

% ------------------------------------------------------------
% Formato: A5, tascabile
% ------------------------------------------------------------
\usepackage{geometry}
\geometry{
  paperwidth=148mm,
  paperheight=210mm,
  top=16mm,
  bottom=18mm,
  inner=16mm,
  outer=13mm,
  headheight=14pt,
  headsep=8mm,
  footskip=12mm
}

% ------------------------------------------------------------
% Font: coppia pensata per la lettura prolungata su formato piccolo.
% TeX Gyre Pagella (clone della Palatino) come testo: occhio alto,
% pensato per restare leggibile a corpi piccoli, più caldo della
% Termes usata nella versione A4. TeX Gyre Heros per i titoli, come
% contrappunto lineare. Entrambe fanno parte della collezione
% TeX Gyre, inclusa in qualsiasi installazione TeX Live standard:
% nessun font da installare a parte.
% ------------------------------------------------------------
\usepackage{fontspec}
\setmainfont{TeX Gyre Pagella}[
  Numbers = OldStyle,
  Ligatures = TeX
]
\setsansfont{TeX Gyre Heros}
\setmonofont{TeX Gyre Cursor}[Scale=0.9]

\usepackage{polyglossia}
\setmainlanguage{italian}

\usepackage{amsmath, amssymb, mathtools}
\usepackage{siunitx}
\usepackage{graphicx}
\usepackage{xcolor}
\usepackage{booktabs}
\usepackage{longtable}
\usepackage{microtype}
% Tolleranza più ampia per la giustificazione: utile su una colonna
% stretta come quella di un formato A5, per assorbire le righe che
% altrimenti sporgerebbero leggermente oltre il margine.
\emergencystretch=3em
\tolerance=1000
\hbadness=3000
\usepackage{setspace}
\onehalfspacing

% Colore discreto per titoli e dettagli grafici (bruno-ocra, richiamo
% cromatico alla roccia e ai pigmenti citati nel testo, non decorativo)
\definecolor{ocra}{RGB}{123,62,32}

% ------------------------------------------------------------
% Intestazioni di pagina: nome parte sulle pagine pari,
% titolo capitolo sulle pagine dispari, numero di pagina in basso
% ------------------------------------------------------------
\usepackage{fancyhdr}
\pagestyle{fancy}
\fancyhf{}
\renewcommand{\headrulewidth}{0.4pt}
\renewcommand{\headrule}{\color{ocra}\hrule height\headrulewidth width\headwidth}
\fancyhead[LE]{\small\itshape\leftmark}
\fancyhead[RO]{\small\itshape\rightmark}
\fancyfoot[LE,RO]{\small\thepage}

\fancypagestyle{plain}{%
  \fancyhf{}
  \fancyfoot[LE,RO]{\small\thepage}
  \renewcommand{\headrulewidth}{0pt}
}

% Marcatori usati da fancyhdr: nome della parte a sinistra (pagine pari),
% titolo del capitolo a destra (pagine dispari). La classe book non
% definisce di suo un comando per marcare le parti nell'intestazione:
% lo si aggancia qui direttamente al comando \part.
\let\oldpart\part
\renewcommand{\part}[1]{\oldpart{#1}\markboth{#1}{}}
\renewcommand{\chaptermark}[1]{\markright{#1}{}}

% ------------------------------------------------------------
% Bibliografia
% ------------------------------------------------------------
\usepackage[sort&compress]{natbib}
\bibliographystyle{plainnat}
% Il testo del libro non contiene citazioni puntuali (\citep, \citet):
% e' un'opera divulgativa in cui le fonti sono richiamate discorsivamente
% nel testo, non con rimandi numerici. Percio' bibtex, senza istruzioni
% aggiuntive, non stamperebbe alcuna voce: nessuna fonte risulterebbe
% "citata" nel senso tecnico del comando \cite. \nocite{*} forza la
% stampa dell'intera bibliografia in references.bib, cosi' come una
% sezione "Fonti" a fine libro elenca tutto il materiale consultato,
% indipendentemente da singoli richiami puntuali nel testo.
\usepackage{etoolbox}
\AtBeginDocument{\nocite{*}}

\usepackage[colorlinks=true,
            linkcolor=ocra,
            citecolor=ocra,
            urlcolor=ocra,
            filecolor=ocra]{hyperref}

% Glossario semplice: comando per termini tecnici con rimando al glossario finale
\newcommand{\termine}[1]{\textit{#1}}

% ------------------------------------------------------------
% Stile dei capitoli e delle parti
% ------------------------------------------------------------
\usepackage{titlesec}

\titleformat{\chapter}[display]
  {\normalfont\sffamily\Large\bfseries\color{ocra}}
  {\chaptertitlename\ \thechapter}
  {6pt}
  {\normalfont\sffamily\LARGE\bfseries\color{black}}
\titlespacing*{\chapter}{0pt}{10pt}{24pt}

\titleformat{\part}[display]
  {\normalfont\sffamily\huge\bfseries\color{ocra}\centering}
  {\partname\ \thepart}
  {12pt}
  {\normalfont\sffamily\Huge\bfseries\color{black}\centering}

\titleformat{\section}
  {\normalfont\sffamily\large\bfseries}
  {\thesection}{0.8em}{}

\titleformat{\subsection}
  {\normalfont\sffamily\normalsize\bfseries}
  {\thesubsection}{0.8em}{}

% ------------------------------------------------------------
% Capolettera per l'apertura di ogni capitolo (tocco tipografico
% da libro tascabile curato, non un semplice muro di testo)
% ------------------------------------------------------------
\usepackage{tikz}
\usepackage{lettrine}
\setcounter{DefaultLines}{3}
\renewcommand{\DefaultLoversize}{0.1}
\renewcommand{\DefaultLraise}{0.0}
\newcommand{\aprecapitolo}[1]{\lettrine[lines=3, findent=2pt, nindent=0pt]{\color{ocra}#1}{}}

% Evita la lettera vietata "—": non caricare pacchetti che la introducano
% automaticamente (nessuna azione necessaria, promemoria per chi edita
% il sorgente a mano)
